% ============================================================================
% Causally Calibrated Model-Based Reinforcement Learning for
% Farmland Consolidation via Learned Environment Dynamics
% ============================================================================
% Target: Computers, Environment and Urban Systems (CEUS)
% ============================================================================
\documentclass[11pt]{article}

\usepackage{amsmath,amssymb}
\usepackage{booktabs}
\usepackage{graphicx}
\usepackage{multirow}
\usepackage{caption}
\usepackage{subcaption}
\usepackage{hyperref}
\usepackage{color}
\usepackage{natbib}
\bibliographystyle{plainnat}

\begin{document}

\title{A Causally Calibrated Learned Environment for Reinforcement Learning-Based Farmland Consolidation Planning}

\author{Anonymous authors}

\date{}
\maketitle

\begin{abstract}
Deep reinforcement learning (DRL) for spatial planning can be limited by expensive environment interactions: county-level farmland consolidation optimization requires 8--12 hours on A100 GPUs per training run, and some runs fail because policies drift toward reward-model exploitation. We propose a model-based geocomputational framework that replaces parcel-level simulation with a lightweight learned environment for reinforcement-learning-based farmland consolidation planning. A 237K-parameter transition model learns state dynamics from 12{,}000 trajectories collected under random and greedy behavioural policies and supports policy training on CPU. To reduce reward exploitation in the learned environment, we introduce causal reward calibration, which uses propensity-score-matched average treatment effect estimates from real trajectory data to rescale the reward differential between high- and low-quality actions. Experiments across 13 townships, 2{,}600 consolidation blocks and 52{,}515 parcels show that uncalibrated model-based policies achieve $-0.976\% \pm 0.129\%$ slope reduction on the real environment; GPU-trained model-free baselines evaluated under the same county-level task achieved $-0.84\%$ and $-0.79\%$. Causal calibration improves performance to $-1.102\% \pm 0.100\%$, a statistically significant 13.0\% gain over the uncalibrated learned environment ($p = 0.004$), and the causally derived calibration factor ($\alpha = 0.185$) is within 3\% of the grid-search optimum. The complete pipeline runs without GPU support in under 2 hours, providing a reproducible decision-support route for large-scale land-resource planning.
\end{abstract}

\noindent\textbf{Keywords:} geocomputation; model-based reinforcement learning; learned environment; causal reward calibration; farmland consolidation; spatial planning; land-resource management


% ============================================================================
\section{Introduction}
\label{sec:intro}
% ============================================================================

Reinforcement learning for spatial planning problems faces a fundamental computational bottleneck: the environment is expensive to simulate. In farmland consolidation optimization, a planning agent allocates a limited exchange budget across thousands of candidate intervention blocks to reduce farmland slope while preserving contiguous consolidated fields. Each environment step updates parcel-level land-use states, recomputes adjacency-dependent farmland contiguity, and evaluates connected-component-based policy targets such as \textit{baimu fang} (100-mu consolidated fields, $\geq$6.67\,ha). In the county-scale environment used in this study, a model-free policy trained for 500{,}000 timesteps requires 8--12 hours on NVIDIA A100 GPUs. Some runs also exhibit \textit{policy drift}: the agent learns to maximize secondary consolidated-field rewards while sacrificing the primary slope-reduction objective.

Model-based reinforcement learning (MBRL) offers a principled solution: learn a \textit{transition model} that approximates the environment's state dynamics, then train the policy on the learned model instead of the real environment. The Dreamer family \citep{hafner2023dreamerv3} demonstrated this approach in continuous control, achieving human-level performance across diverse domains by ``dreaming'' trajectories in a learned latent space. However, MBRL has not been applied to large-scale spatial planning problems, where the state space (44{,}212 dimensions for county-level farmland optimization) and the discrete, combinatorial action space (2{,}600 blocks) pose distinct challenges.

This paper makes three contributions:

\begin{enumerate}
\item \textbf{A learned environment for spatial planning.} We train a neural transition model (237K parameters) that predicts next-step observations and rewards from current state and action, achieving 0.9998 cosine similarity with ground-truth transitions. The model trains in 60 minutes on CPU from 12{,}000 trajectory samples.

\item \textbf{Causal reward calibration.} We introduce a mechanism that uses propensity-score-matched average treatment effect (ATT) estimates from real trajectory data to calibrate the learned reward function. The calibration corrects a 5.4$\times$ overestimation of action quality differentials, reducing reward exploitation and improving downstream policy quality by 13.0\% ($p = 0.004$). A grid search over 8 reward scaling factors confirms that the causally derived $\alpha = 0.185$ falls within 3\% of the empirical optimum, supporting causal inference as a data-driven calibration method in this setting.

\item \textbf{Order-of-magnitude compute reduction.} The complete pipeline---trajectory collection, transition model training, and policy optimization---requires no GPU and completes in under 2 hours, compared to 8--12 GPU-hours for model-free training. Model-based policies evaluated on the real environment achieve $-0.976\%$ slope reduction (without calibration) or $-1.102\%$ (with calibration). For context, the model-free baselines achieve $-0.84\%$ for MARL and $-0.79\%$ for centralized DRL; the formal statistical test in this paper is the 15-seed comparison between calibrated and uncalibrated learned-environment policies.
\end{enumerate}


% ============================================================================
\section{Related Work}
\label{sec:related}
% ============================================================================

\subsection{Model-based reinforcement learning}

MBRL methods learn environment dynamics to reduce sample complexity. World Models \citep{ha2018world} introduced the concept of training policies in ``dreams''---trajectories generated by a learned latent dynamics model. Dreamer V3 \citep{hafner2023dreamerv3} extended this to diverse continuous control domains using Recurrent State-Space Models. MuZero \citep{schrittwieser2020muzero} achieved superhuman performance in board games and Atari by planning with a learned model. In spatial domains, \citet{zheng2023spatial} applied model-free DRL to urban community planning, but applications of MBRL to large-scale spatial planning remain limited.

Our approach differs from Dreamer-style methods in a key respect: rather than jointly learning a latent state representation and dynamics model, we operate directly in the environment's native observation space. This is feasible because the farmland optimization environment already provides a compact, informative state representation (17 features per block $\times$ 2{,}600 blocks + 12 global features), eliminating the need for learned encoding.

\subsection{Reward model exploitation}

The problem of agents exploiting learned reward models is well-documented. \citet{amodei2016concrete} identified reward hacking as a key safety concern. In offline RL, \citet{levine2020offline} showed that out-of-distribution actions can exploit value function overestimation. Our causal calibration approach is related to reward model correction methods in RLHF \citep{bai2022training}, but uses causal inference rather than human feedback to ground the correction.

\subsection{DRL for farmland consolidation}

Farmland consolidation is a spatial resource-allocation problem with terrain, land-use, and contiguity constraints. Direct parcel-level decisions are difficult because cadastral parcels are numerous, irregular, and linked through adjacency relationships. This study therefore uses connected blocks of swappable farmland and forest parcels as planning units. The block abstraction keeps the action space tractable while retaining the parcel-level information needed to evaluate slope change, farmland contiguity, and large contiguous farmland patches. The model-free baselines in this paper use the same block-level county environment defined below, so the learned-environment results can be evaluated against policies trained directly on the real parcel-simulation environment.

\subsection{Causal inference for reward calibration}

Propensity score methods \citep{stuart2010matching} enable treatment effect estimation from observational data by balancing confounders between treated and control groups. We adapt this framework to the RL setting: ``treatment'' is the selection of a high-potential block, ``outcome'' is the resulting reward, and ``confounders'' are the global state features that influence both block selection and reward. This connection provides a data-driven way to estimate how strongly the learned environment overstates action-quality differences, without treating reward scaling as a manually tuned hyperparameter.


% ============================================================================
\section{Problem Setting}
\label{sec:problem}
% ============================================================================

The study area contains all 13 townships of Bishan District, Chongqing, southwest China. The county-level dataset contains 52{,}515 swappable farmland and forest parcels aggregated into 2{,}600 candidate consolidation blocks. A parcel $i$ has a land-use state $y_i \in \{\text{farmland}, \text{forest}\}$, area $a_i$, slope $S_i$, and adjacency set $\mathcal{N}(i)$. A block $b$ is valid at timestep $t$ only if it still contains at least one unswapped farmland parcel and at least one unswapped forest parcel. This validity condition defines the action mask used by all DRL policies.

Selecting a block triggers a local paired-swap transition rather than an arbitrary land-use change. For up to $\kappa=5$ paired swaps, the environment selects a high-slope farmland parcel for conversion to forest and a low-slope forest parcel for conversion to farmland within the chosen block. A pair is executed only when it reduces local slope, after which parcel land-use states, farmland-neighbor counts, block features, and county-level metrics are updated. This operator preserves the total number of swappable parcels while moving farmland toward lower-slope and more connected configurations.

\textit{Baimu fang} metrics are computed from the county-wide farmland adjacency graph. A union--find procedure identifies connected farmland components, sums their areas, and counts a component as \textit{baimu fang} if its area is at least 66{,}700\,m$^2$ (6.67\,ha). Because adjacency includes parcel edges across township borders, a consolidated farmland patch may span administrative boundaries.

The centralized state $\mathbf{s}_t \in \mathbb{R}^{D}$ ($D = 44{,}212$) concatenates per-block features for all 2{,}600 blocks ($\mathbf{X}_t \in \mathbb{R}^{2600 \times 17}$) and 12 county-level global metrics ($\mathbf{g}_t \in \mathbb{R}^{12}$). The 17 block features include farmland and forest slope summaries, best available slope gain, available farmland and forest area, swap-potential ratio, compactness, neighboring-block investment context, current block farmland area, and an invested-block indicator. The global features encode remaining budget, current slope and contiguity, episode progress, slope and contiguity improvement, \textit{baimu fang} count and area fraction, invested-block fraction, investment entropy across townships, and maximum single-township investment share.

The action space is $\mathcal{A}=\{1,\ldots,2600\}$ with invalid blocks masked. Each episode has a total budget of 500 paired swaps and uses $\kappa=5$ swaps per selected block, yielding 100 decision steps. The reward combines county-wide slope reduction, contiguity change, \textit{baimu fang} area change, and \textit{baimu fang} count change. This multi-objective reward is intentionally retained across all methods so that model-based and model-free policies face the same planning trade-off.

Two model-free baselines are trained directly on this real environment. The centralized DRL baseline uses Maskable PPO over the full 2{,}600-block action space and the 44{,}212-dimensional state. The MARL baseline decomposes the county into 13 township agents that share one block-scoring policy; agents act in turn, observe local padded block features plus global metrics, and receive the same county-level reward. Both model-free baselines use 500{,}000 training timesteps on A100 GPUs. The centralized baseline takes approximately 8 hours per seed, and the MARL baseline takes approximately 12 hours per seed. Across 9 model-free seeds, 3 showed severe policy drift, defined as more than 50\% positive-slope episodes during training.


% ============================================================================
\section{Methods}
\label{sec:methods}
% ============================================================================

Our approach consists of three stages (Fig.~\ref{fig:pipeline}): (1) trajectory collection from the real environment under diverse behavioural policies, (2) transition model training on collected trajectories, and (3) policy optimization on the learned environment with optional causal reward calibration.

\begin{figure}[htbp]
\centering
\fbox{\begin{minipage}{0.92\textwidth}
\centering
Real county environment $\rightarrow$ random and greedy trajectory collection $\rightarrow$ neural transition model $\rightarrow$ learned county environment $\rightarrow$ policy training with optional causal reward calibration $\rightarrow$ deterministic evaluation on the real county environment.
\end{minipage}}
\caption{Pipeline for learned-environment policy training. The real parcel-simulation environment is used to collect trajectories and to evaluate final policies. The learned environment is used only for policy optimization, with causal reward calibration applied during model-based training.}
\label{fig:pipeline}
\end{figure}

\subsection{Stage 1: Trajectory collection}

We collect state-action-reward-next\_state tuples $\{(\mathbf{s}_t, a_t, r_t, \mathbf{s}_{t+1})\}$ from the real CountyLevelEnv using two behavioural policies:

\begin{itemize}
\item \textbf{Random}: uniform random selection from valid (unmasked) blocks. 3 seeds $\times$ 20 episodes = 6{,}000 transitions.
\item \textbf{Greedy}: select the block with highest slope reduction potential (feature index 3: best\_swap\_gain\_norm). 3 seeds $\times$ 20 episodes = 6{,}000 transitions.
\end{itemize}

The combined dataset of 12{,}000 transitions provides both exploratory coverage (random) and high-quality demonstrations (greedy). Each transition stores structured block-level features ($\mathbb{R}^{2600 \times 17}$, float16) and global features ($\mathbb{R}^{12}$, float32), totalling 1.6\,GB compressed.

\subsection{Stage 2: Neural transition model}

The \textit{TransitionModel} predicts the next observation and reward given current state and action:
\begin{equation}
\hat{\mathbf{s}}_{t+1}, \hat{r}_t = f_\theta(\mathbf{s}_t, a_t)
\end{equation}

\textbf{Architecture.} The model consists of four components:
\begin{enumerate}
\item \textit{Block encoder}: a shared MLP ($17 \to 64 \to 32$) that encodes each block's features independently, producing $\mathbf{h}_b \in \mathbb{R}^{32}$ for each of the 2{,}600 blocks.
\item \textit{Action embedding}: $\text{Embed}(a_t) \in \mathbb{R}^{32}$, a learned embedding for the selected block.
\item \textit{Global encoder}: an MLP ($12 \to 64 \to 32$) encoding the global state.
\item \textit{Context aggregation}: the encoding of the selected block, the action embedding, the global encoding, and the mean-pooled block encodings (spatial context) are concatenated into a 128-dimensional context vector.
\end{enumerate}

From this context, three heads predict:
\begin{itemize}
\item $\Delta\mathbf{x}_{a_t} \in \mathbb{R}^{17}$: the residual change to the \textit{selected block's} features (via MLP $128 \to 256 \to 256 \to 17$). All other blocks remain unchanged: $\hat{\mathbf{x}}_{b,t+1} = \mathbf{x}_{b,t}$ for $b \neq a_t$.
\item $\Delta\mathbf{g} \in \mathbb{R}^{12}$: the residual change to global features (via MLP $128 \to 256 \to 12$).
\item $\hat{r}_t \in \mathbb{R}$: the predicted reward (via MLP $128 \to 64 \to 1$).
\end{itemize}

The residual formulation is critical: inter-step state changes are small (most blocks are unaffected by any single action), so predicting $\Delta$ rather than absolute values avoids the identity-mapping learning problem. Total parameters: 236{,}958.

\textbf{Training.} The model is trained with Adam ($\text{lr} = 10^{-3}$, weight decay $10^{-5}$) for 50 epochs on the 12{,}000-transition dataset (90/10 train/val split). Loss: $\mathcal{L} = \text{MSE}(\hat{\mathbf{X}}_{t+1}, \mathbf{X}_{t+1}) + \text{MSE}(\hat{\mathbf{g}}_{t+1}, \mathbf{g}_{t+1}) + 0.1 \cdot \text{MSE}(\hat{r}_t, r_t)$. Training completes in $\sim$60 minutes on CPU. Final validation cosine similarity: 0.9998.

\subsection{Stage 3: Policy optimization on learned environment}

The \textit{LearnedCountyEnv} wraps the trained transition model as a Gymnasium-compatible environment with the same observation space, action space, and action masking interface as the real CountyLevelEnv. At each step, $\texttt{env.step}(a)$ calls the transition model's forward pass instead of executing real parcel-level swaps. The initial state is drawn from the trajectory dataset.

Policy optimization uses Maskable PPO with a block-scoring actor--critic policy. The actor scores each candidate block from its 17 block features and the 12 global features, masks invalid blocks, and samples from the remaining block logits. The value network estimates the county-level state value from the global features. Model-based training uses lr $= 3 \times 10^{-4}$, $n_\text{steps} = 256$, $\gamma = 0.995$, batch size 128, 10 PPO epochs, and entropy coefficient 0.005. Training for 100{,}000 timesteps produces 1{,}000 episodes and completes in 25--57 minutes on CPU.

\textbf{Critical distinction}: the policy is trained on the learned environment but \textit{evaluated on the real environment}. All reported metrics (slope change, contiguity, \textit{baimu fang}) are measured on the real CountyLevelEnv with deterministic inference.

\subsection{Causal reward calibration}
\label{sec:calibration}

A key failure mode of model-based RL is \textit{reward exploitation}: the policy learns to select actions that receive high reward from the learned model but perform poorly in the real environment. We observe this as systematic overestimation of the reward differential between ``good'' and ``bad'' actions---the learned model amplifies the quality gap by 5.4$\times$ compared to empirical data.

To correct this, we estimate the true causal effect of action quality on reward using propensity score stratification on the trajectory dataset:

\textbf{Treatment definition.} For each transition, the ``treatment'' is binary: $T_t = 1$ if the selected block's slope gap (feature 3) exceeds the median across all blocks, $T_t = 0$ otherwise.

\textbf{Confounders.} Global state features (budget remaining, current slope, contiguity, step fraction, slope improvement) and selected block features (farmland slope, forest slope, swap potential, investment status) serve as confounders that influence both treatment assignment and outcome.

\textbf{ATT estimation.} Propensity scores are estimated via gradient-boosted trees with 5-fold cross-validation. After trimming extreme scores ($\hat{e} \notin [0.05, 0.95]$), the ATT is computed via stratified matching across 5 propensity score strata with bootstrap standard errors (200 resamples). Covariate balance after matching was verified via standardized mean differences (all $< 0.1$; balance diagnostics available in supplementary materials).

\textbf{Calibration.} Let $\hat{\tau}_\text{ATT}$ be the empirical treatment effect and $\hat{\tau}_\text{model}$ be the learned model's predicted treatment effect (mean reward for high-potential blocks minus low-potential blocks). The calibration factor is:
\begin{equation}
\alpha = \text{clip}\left(\frac{\hat{\tau}_\text{ATT}}{\hat{\tau}_\text{model}},\; 0.1,\; 5.0\right)
\end{equation}

During model-based policy training, all rewards from the learned environment are scaled by $\alpha$. In our experiments, $\hat{\tau}_\text{ATT} = +0.049$, $\hat{\tau}_\text{model} = +0.265$, yielding $\alpha = 0.185$---the learned model overestimates the benefit of good actions by 5.4$\times$, and the calibration shrinks rewards accordingly.


% ============================================================================
\section{Experiments}
\label{sec:experiments}
% ============================================================================

\subsection{Experimental setup}

All experiments use the same county-level farmland consolidation environment defined in Section~\ref{sec:problem}: 13 Bishan District townships, 2{,}600 blocks, 52{,}515 parcels, and a budget of 500 paired swaps. Model-based policies are trained in the learned environment and then evaluated in the real environment. Model-free baselines are trained and evaluated directly in the real environment under the same budget, reward, action mask, and deterministic evaluation protocol.

\textbf{Methods compared:}
\begin{itemize}
\item \textbf{Greedy-Sequential}: deterministic baseline that visits valid blocks sorted by initial farmland slope, investing one step per block before cycling.
\item \textbf{Centralized DRL} (model-free): Maskable PPO on the real county environment, 5 seeds, 500K steps, A100 GPU.
\item \textbf{MARL DRL} (model-free): 13-agent turn-based shared-policy DRL on the real county environment, 4 seeds, 500K steps, A100 GPU.
\item \textbf{Model-Based} (ours): Maskable PPO on LearnedCountyEnv, 15 seeds, 100K steps, CPU.
\item \textbf{Model-Based + Calibration} (ours): same with causal reward calibration ($\alpha = 0.185$), 15 seeds.
\item \textbf{Heuristic Scaling} (ablation): reward scaled by the values tested in Section~\ref{sec:grid_search}, 5 seeds each.
\end{itemize}

\subsection{Main results}

Table~\ref{tab:main} presents the main comparison.

\begin{table}[htbp]
\centering
\caption{Method comparison. All slope/contiguity metrics are evaluated on the real county environment. Training time includes all pipeline stages (trajectory collection + model training + policy optimization) for model-based methods. Asterisks indicate Mann-Whitney $U$ tests against Greedy-Sequential; comparisons with model-free DRL are reported descriptively because seed counts and training architectures differ. Best result in bold.}
\label{tab:main}
\begin{tabular}{lcccccc}
\toprule
Method & Seeds & Slope (\%) & Sig. & Cont. & Time & Hardware \\
\midrule
Greedy-Sequential      & 1  & $-0.27$              & ---  & $+0.008$ & $10$\,s   & CPU \\
Centralized DRL        & 5  & $-0.79 \pm 0.36$     &      & $+0.017$ & $8$\,h    & A100 \\
MARL DRL               & 4  & $-0.84 \pm 0.07$     &      & $+0.017$ & $12$\,h   & A100 \\
\midrule
Model-Based (ours)     & 15 & $-0.976 \pm 0.129$   & $**$ & $+0.013$ & $\sim$2\,h & CPU \\
\textbf{MB + Calibration} & 15 & $\mathbf{-1.102 \pm 0.100}$ & $**$ & $+0.011$ & $\sim$2\,h & CPU \\
\bottomrule
\end{tabular}
\end{table}

\textbf{Model-based RL improves mean slope reduction relative to model-free baselines.} Without calibration, model-based policies achieve $-0.976\%$ slope reduction across 15 seeds, compared with MARL ($-0.84\%$) and centralized DRL ($-0.79\%$). With causal calibration, performance improves to $-1.102\%$. The strongest formal test in this experiment is the paired methodological comparison between calibrated and uncalibrated learned environments: causal calibration improves model-based policy quality by 13.0\% ($p = 0.004$, Mann-Whitney $U$ test). The comparison with model-free baselines should be read as an external baseline comparison under the same task rather than as a fully balanced statistical test.

\textbf{Compute reduction.} The model-based pipeline completes in $\sim$2 hours on CPU (trajectory collection: $\sim$40\,min; transition model: $\sim$60\,min; policy: $\sim$45\,min), compared to 8--12 GPU-hours for model-free methods. This represents a $\sim$16$\times$ wall-clock speedup with no GPU requirement.

\textbf{Improved reliability.} Model-based training exhibits no policy drift: all 15 seeds converge to consistent slope reduction ($-0.73\%$ to $-1.32\%$, std $= 0.129\%$). In contrast, model-free training shows 33\% failure rate (3/9 seeds across centralized and MARL architectures).

\subsection{Transition model quality}
\label{sec:tm_quality}

\begin{table}[htbp]
\centering
\caption{Transition model validation metrics.}
\label{tab:tm}
\begin{tabular}{lcc}
\toprule
Metric & Value & Threshold \\
\midrule
Cosine similarity (obs) & 0.9998 & $>$0.95 \\
Best validation loss & 0.088 & --- \\
Reward MSE & 1.097 & --- \\
Parameters & 236{,}958 & --- \\
Training time & 60\,min & --- \\
Training data & 12{,}000 transitions & --- \\
\bottomrule
\end{tabular}
\end{table}

The transition model achieves high aggregate state-prediction similarity (cosine similarity 0.9998). This metric is useful as a sanity check but should be interpreted cautiously because most blocks do not change after any single action. The reward MSE is higher (1.097), reflecting the high variance of rewards across random and greedy policies. The causal calibration mechanism (Section~\ref{sec:calibration}) therefore targets systematic reward bias rather than claiming perfect per-step reward prediction.

\subsection{Ablation: Causal reward calibration (E5)}

\begin{table}[htbp]
\centering
\caption{Effect of causal reward calibration on policy quality (15 seeds). All metrics evaluated on real env.}
\label{tab:e5}
\begin{tabular}{lcccc}
\toprule
Configuration & $n$ & Slope (\%) & $p$-value & Wins \\
\midrule
No calibration ($\alpha = 1.0$)       & 15 & $-0.976 \pm 0.129$ & --- & --- \\
With calibration ($\alpha = 0.185$)   & 15 & $\mathbf{-1.102 \pm 0.100}$ & $0.004^{**}$ & 10/15 \\
\midrule
Improvement & & $+13.0\%$ & & \\
\bottomrule
\end{tabular}
\end{table}

Causal calibration improves slope reduction by 13.0\% ($-1.102\%$ vs.\ $-0.976\%$), statistically significant at $p = 0.004$ (Mann-Whitney $U = 48$, one-sided). The calibrated policy achieves better slope reduction than the uncalibrated version in 10 of 15 seeds. The improvement is robust: the worst calibrated seed ($-0.946\%$) still matches the uncalibrated mean.

The calibration factor $\alpha = 0.185$ indicates that the learned environment overestimates the reward differential between high- and low-quality actions by $5.4\times$. Without correction, the policy exploits this overestimation by concentrating on a narrow set of ``apparently high-value'' blocks; with calibration, the flattened reward landscape encourages broader exploration and more balanced block selection.

\subsection{Causal calibration vs.\ heuristic reward scaling}
\label{sec:grid_search}

A natural question is whether the benefit of causal calibration could be replicated by simply tuning a reward scaling hyperparameter via grid search. To test this, we trained policies with 8 reward scaling factors, from 0.1 to 1.0 and including the causally derived value 0.185, using 5 seeds per value (Table~\ref{tab:grid}).

\begin{table}[htbp]
\centering
\caption{Reward scaling grid search (5 seeds each). The causally derived $\alpha = 0.185$ is marked. Grid-search optimum at $\alpha = 0.200$.}
\label{tab:grid}
\begin{tabular}{lcc}
\toprule
$\alpha$ & Slope (\%) & Std \\
\midrule
$0.100$                          & $-1.123$ & $0.151$ \\
$0.150$                          & $-1.120$ & $0.109$ \\
$0.185$ \textit{(causal)}        & $-1.171$ & $0.103$ \\
$\mathbf{0.200}$ \textit{(grid best)} & $\mathbf{-1.209}$ & $\mathbf{0.092}$ \\
$0.300$                          & $-1.171$ & $0.077$ \\
$0.500$                          & $-1.043$ & $0.065$ \\
$0.700$                          & $-0.991$ & $0.047$ \\
$1.000$ (no scaling)             & $-0.959$ & $0.055$ \\
\bottomrule
\end{tabular}
\end{table}

The results reveal three key findings. First, the optimal scaling region is $\alpha \in [0.15, 0.30]$, with performance degrading monotonically for $\alpha > 0.3$. Second, the causally derived $\alpha = 0.185$ achieves $-1.171\%$, within 0.038 percentage points (3.1\%) of the grid-search optimum $\alpha = 0.200$ ($-1.209\%$). Third, the no-scaling baseline ($\alpha = 1.0$, $-0.959\%$) is the weakest performer, supporting the interpretation that reward overestimation affects policy training in this setting.

These results support causal calibration as a data-driven scaling method: the causally derived $\alpha$ lands in the optimal plateau without any hyperparameter search, which would otherwise require $8 \times 5 = 40$ additional training runs. We note that a state-dependent calibration factor $\alpha(\mathbf{s})$ could potentially improve upon the global constant; we leave this extension to future work.

\subsection{Ablation: GeoFM embedding augmentation (E4)}

We tested augmenting the transition model's per-block input with 64-dimensional AlphaEarth GeoFM embeddings \citep{brown2025alphaearth}, extracted at each block's centroid via Google Earth Engine.

\begin{table}[htbp]
\centering
\caption{GeoFM embedding ablation. TM = transition model.}
\label{tab:e4}
\begin{tabular}{lcccc}
\toprule
Configuration & TM Params & TM Val Loss & Policy Slope (\%) \\
\midrule
No GeoFM (17-dim input)      & 236{,}958 & 0.085 & $-0.96 \pm 0.14\%$ \\
With GeoFM (81-dim input)    & 241{,}054 & \textbf{0.051} & $-0.70 \pm 0.01\%$ \\
\bottomrule
\end{tabular}
\end{table}

GeoFM embeddings reduce transition model validation loss by 40\% (0.051 vs.\ 0.085), indicating better one-step state prediction. However, downstream policy quality \textit{decreases} ($-0.70\%$ vs.\ $-0.96\%$). This counterintuitive result suggests that a more accurate one-step model is not necessarily a better training environment for policy optimization. The finding aligns with observations in the MBRL literature that model accuracy and policy quality are not monotonically related \citep{lambert2022investigating}.


% ============================================================================
\section{Discussion}
\label{sec:discussion}
% ============================================================================

\subsection{Why model-based training can perform well with an approximate environment}

The observed advantage of model-based policies in this benchmark is initially counterintuitive: policies trained on an approximate model achieved lower mean slope than policies trained directly on the real environment, although the comparison with model-free baselines remains descriptive because seed counts and training architectures differ. Three factors may contribute:

\textbf{Noise-free gradients.} The learned environment provides smooth, deterministic state transitions (given a fixed transition model), eliminating the stochastic variation present in the real environment. This produces cleaner policy gradients and more stable optimization.

\textbf{Faster iteration.} At 100{,}000 timesteps, the model-based agent sees 1{,}000 full episodes, while the model-free agent at the same step count has seen the same number. However, the model-based steps are $\sim$$100\times$ faster (no parcel-level simulation), allowing more gradient updates per wall-clock minute.

\textbf{Implicit curriculum from trajectory distribution.} The training data includes both random (exploratory) and greedy (exploitative) trajectories. The transition model learns an ``average'' dynamics that smooths over the stochastic greedy engine's variability, providing a more learnable training signal than the real environment's episodic noise.

\subsection{Causal calibration as reward regularization}

The calibration factor $\alpha = 0.185$ effectively shrinks the reward scale by $5.4\times$, which has a regularization effect similar to reward clipping or normalization. However, unlike generic normalization, causal calibration is \textit{grounded in empirical evidence}: the scaling is derived from the observed treatment effect of action quality on reward, not from arbitrary hyperparameter tuning. The grid search experiment (Section~\ref{sec:grid_search}) confirms this: the causally derived $\alpha = 0.185$ falls within the optimal plateau $[0.15, 0.30]$, matching the grid-search optimum ($\alpha = 0.200$) to within 3\%. This suggests a possible route for transfer: in a new environment, the same estimation procedure could recompute $\alpha$ from local trajectories instead of reusing a fixed hand-tuned value.

\subsection{The accuracy--utility paradox in GeoFM augmentation}

The GeoFM ablation reveals a paradox: adding rich spatial context (64-dimensional foundation model embeddings) improves transition model accuracy but hurts policy quality. We hypothesize two mechanisms:

\textbf{Over-faithful modeling.} A more accurate model more faithfully reproduces the real environment's challenges---including the reward landscape features that cause policy drift in model-free training. A less accurate model may accidentally ``smooth over'' these degenerate attractors, making the optimization landscape more benign.

\textbf{Static feature contamination.} GeoFM embeddings are static (fixed for the 2020 snapshot) while block features change every step. The constant presence of 64 high-dimensional features may dominate the block encoder's representation, masking the dynamic features that the policy needs to respond to.

\subsection{Limitations}

\textbf{Distribution shift.} The transition model is trained on random and greedy trajectories. If the learned policy visits states far from this training distribution, prediction quality may degrade. Dyna-style hybrid training (alternating real and learned steps) could mitigate this but was not explored here.

\textbf{Single study area.} All experiments use Bishan District. Cross-county generalization requires new trajectory data, though the transition model architecture is region-agnostic.

\textbf{Global calibration factor.} The current calibration factor $\alpha$ is a global constant. A state-dependent $\alpha(\mathbf{s})$ could capture heterogeneous reward overestimation across different planning stages, potentially improving late-episode performance where distribution shift is most pronounced.

\textbf{Reward prediction.} While state prediction is near-perfect (cosine 0.9998), reward prediction has higher error (MSE 1.097). The causal calibration addresses systematic bias but not per-step noise.


% ============================================================================
\section{Conclusions}
\label{sec:conclusions}
% ============================================================================

We have demonstrated that model-based reinforcement learning, combined with causal reward calibration, can reduce computational cost and, in this benchmark, improve mean slope reduction for county-level farmland consolidation planning. The key findings are:

\begin{enumerate}
\item \textbf{Learned environments are viable for large-scale spatial planning.} A 237K-parameter transition model trained on 12{,}000 trajectory samples achieves 0.9998 cosine similarity with real environment dynamics and supports effective policy training.

\item \textbf{Model-based policies improve mean slope reduction relative to model-free baselines.} Policies trained on the learned environment achieve $-0.976\%$ slope reduction on the real environment across 15 seeds, compared with 12-hour A100-trained MARL ($-0.84\%$) and centralized DRL ($-0.79\%$).

\item \textbf{Causal calibration addresses reward exploitation.} Propensity-score-matched treatment effect estimates reveal that the learned reward function overestimates action quality differentials by $5.4\times$. Correcting this via a calibration factor improves policy quality by 13.0\% to $-1.102\%$ ($p = 0.004$). The causally derived calibration factor matches the grid-search optimum to within 3\%, supporting causal inference as a data-driven calibration method in this setting.

\item \textbf{An order-of-magnitude compute reduction is achievable.} The complete model-based pipeline runs in $\sim$2 hours on CPU, compared to 8--12 hours on A100 for model-free training, with zero GPU requirement.

\item \textbf{Model accuracy and policy quality are not monotonically related.} Adding GeoFM embeddings improves transition model accuracy by 40\% but degrades policy quality, indicating that one-step predictive accuracy alone is not sufficient to judge a learned environment for policy optimization.
\end{enumerate}

These results suggest that the ``dream then act'' paradigm---collecting diverse experience, learning environment dynamics, and optimizing policies in imagination---is a practical route for county-scale farmland consolidation environments where parcel-level simulation is expensive. The causal calibration mechanism provides a principled, data-driven response to the reward exploitation pattern observed in this study.


% ============================================================================
\section*{References}
% ============================================================================

\begin{thebibliography}{25}

\bibitem[Amodei et~al.(2016)]{amodei2016concrete}
Amodei, D., Olah, C., Steinhardt, J., et~al. (2016).
Concrete problems in AI safety. arXiv:1606.06565.

\bibitem[Bai et~al.(2022)]{bai2022training}
Bai, Y., Jones, A., Ndousse, K., et~al. (2022).
Training a helpful and harmless assistant with RLHF. arXiv:2204.05862.

\bibitem[Brown et~al.(2025)]{brown2025alphaearth}
Brown, C., Kazmierski, M., Pasquarella, V., et~al. (2025).
AlphaEarth Foundations: An embedding field model for global mapping. arXiv:2507.22291.

\bibitem[Ha and Schmidhuber(2018)]{ha2018world}
Ha, D. and Schmidhuber, J. (2018).
World Models. arXiv:1803.10122.

\bibitem[Hafner et~al.(2023)]{hafner2023dreamerv3}
Hafner, D., Pasukonis, J., Ba, J., and Lillicrap, T. (2023).
Mastering diverse domains through world models. arXiv:2301.04104.

\bibitem[Lambert et~al.(2022)]{lambert2022investigating}
Lambert, N., Amos, B., Yadan, O., and Calandra, R. (2022).
Objective mismatch in model-based reinforcement learning. In L4DC.

\bibitem[Levine et~al.(2020)]{levine2020offline}
Levine, S., Kumar, A., Tucker, G., and Fu, J. (2020).
Offline reinforcement learning: Tutorial, review, and perspectives. arXiv:2005.01643.

\bibitem[Schrittwieser et~al.(2020)]{schrittwieser2020muzero}
Schrittwieser, J., Antonoglou, I., Hubert, T., et~al. (2020).
Mastering Atari, Go, chess and shogi by planning with a learned model. Nature 588, 604--609.

\bibitem[Stuart(2010)]{stuart2010matching}
Stuart, E.~A. (2010).
Matching methods for causal inference: A review. Statistical Science 25(1), 1--21.

\bibitem[Zheng et~al.(2023)]{zheng2023spatial}
Zheng, Y., Lin, Y., Zhao, L., et~al. (2023).
Spatial planning of urban communities via deep reinforcement learning. Nature Computational Science 3(9), 748--762.

\end{thebibliography}

\end{document}


